\documentclass[11pt,a4paper]{article}

\usepackage{kotex}
\usepackage[margin=25mm]{geometry}
\usepackage{amsmath,amssymb}
\usepackage{booktabs}
\usepackage{xcolor}
\usepackage{listings}
\usepackage{hyperref}
\usepackage{enumitem}

\hypersetup{
  colorlinks=true,
  linkcolor=blue,
  urlcolor=blue
}

\definecolor{codebg}{RGB}{247,247,247}
\definecolor{codeframe}{RGB}{210,210,210}
\definecolor{codecomment}{RGB}{0,110,60}

\lstset{
  language=Python,
  basicstyle=\ttfamily\small,
  backgroundcolor=\color{codebg},
  frame=single,
  rulecolor=\color{codeframe},
  keywordstyle=\color{blue},
  commentstyle=\color{codecomment},
  stringstyle=\color{red!60!black},
  showstringspaces=false,
  breaklines=true,
  columns=fullflexible
}

\title{\textbf{Python for Physicists}\\[0.4em]
\large Chapter 1. Python 기초}
\author{}
\date{}

\begin{document}

\maketitle
\tableofcontents
\newpage

\section{학습 목표}

이번 장을 마치면 다음 내용을 설명하고 사용할 수 있어야 한다.

\begin{itemize}[leftmargin=2em]
  \item \texttt{print()}의 역할을 이해한다.
  \item 변수와 객체의 차이를 설명할 수 있다.
  \item Mutable 객체와 Immutable 객체를 구분할 수 있다.
  \item Python의 기본 자료형을 이해한다.
  \item 부동소수점 오차가 발생하는 이유를 설명할 수 있다.
  \item \texttt{input()}과 형 변환을 사용할 수 있다.
  \item 간단한 BMI 계산기를 작성할 수 있다.
\end{itemize}

\section{출력하기: \texttt{print()}}

Python에서 가장 먼저 접하는 함수는 \texttt{print()}이다.

\begin{lstlisting}
print("Hello, World!")
\end{lstlisting}

실행 결과는 다음과 같다.

\begin{verbatim}
Hello, World!
\end{verbatim}

\texttt{print()}는 계산을 수행하는 함수가 아니라, 결과를 화면에 출력하는 함수이다.

\begin{lstlisting}
3 + 5
\end{lstlisting}

일반적인 Python 스크립트에서는 계산 결과가 자동으로 출력되지 않는다.

\begin{lstlisting}
print(3 + 5)
\end{lstlisting}

\paragraph{핵심}
계산과 출력은 서로 다른 작업이다.

\section{변수와 객체}

\begin{lstlisting}
age = 22
\end{lstlisting}

Python에서 변수는 값을 저장하는 상자라기보다 객체를 가리키는 이름으로 이해하는 편이 정확하다.

\subsection{대입의 의미}

\begin{lstlisting}
age = 22
age = 23
\end{lstlisting}

두 번째 줄은 \texttt{22}를 \texttt{23}으로 수정하는 것이 아니다. 새로운 정수 객체 \texttt{23}을 만든 뒤, \texttt{age}가 그 객체를 가리키도록 바꾸는 것이다.

\subsection{수학적 관점}

변수들의 집합을 \(V\), 객체들의 집합을 \(O\)라고 하면 프로그램의 상태를 다음과 같은 함수로 생각할 수 있다.

\[
f:V\rightarrow O
\]

예를 들어 처음에는

\[
f(\texttt{age})=22
\]

였고, 대입 이후에는

\[
f(\texttt{age})=23
\]

이 된다.

\section{Mutable과 Immutable}

Python 객체는 객체 자체를 수정할 수 있는지에 따라 Mutable과 Immutable로 나눌 수 있다.

\subsection{Immutable 객체}

대표적인 Immutable 자료형은 \texttt{int}, \texttt{float}, \texttt{bool}, \texttt{str}, \texttt{tuple}이다.

\begin{lstlisting}
a = 10
a = 20
\end{lstlisting}

이는 정수 객체 \texttt{10}을 수정하는 것이 아니라, \texttt{a}가 새로운 정수 객체 \texttt{20}을 가리키게 만드는 과정이다.

\subsection{Mutable 객체}

대표적인 Mutable 자료형은 \texttt{list}, \texttt{dict}, \texttt{set}이다.

\begin{lstlisting}
a = [1, 2, 3]
b = a

a.append(4)

print(a)
print(b)
\end{lstlisting}

출력 결과는 다음과 같다.

\begin{verbatim}
[1, 2, 3, 4]
[1, 2, 3, 4]
\end{verbatim}

\texttt{a}와 \texttt{b}가 같은 리스트 객체를 가리키고 있기 때문에, 객체 자체가 수정되면 두 이름을 통해 동일한 변화가 보인다.

\section{기본 자료형}

\begin{center}
\begin{tabular}{ll}
\toprule
자료형 & 설명 \\
\midrule
\texttt{int} & 정수 \\
\texttt{float} & 부동소수점 수 \\
\texttt{str} & 문자열 \\
\texttt{bool} & 참 또는 거짓 \\
\bottomrule
\end{tabular}
\end{center}

\begin{lstlisting}
print(type(10))
print(type(3.14))
print(type("Python"))
print(type(True))
\end{lstlisting}

\section{부동소수점 오차}

\begin{lstlisting}
print(0.1 + 0.2)
\end{lstlisting}

\begin{verbatim}
0.30000000000000004
\end{verbatim}

이것은 Python의 버그가 아니다. 컴퓨터는 유한한 비트만 사용하므로 대부분의 실수를 정확히 표현할 수 없다.

실수 전체의 집합을 \(\mathbb{R}\), 컴퓨터가 표현 가능한 부동소수점 수들의 집합을 \(F\)라고 하면 저장 과정을 다음과 같이 생각할 수 있다.

\[
r:\mathbb{R}\rightarrow F
\]

즉, 실제 값 \(x\) 대신 가장 가까운 표현 가능한 값 \(r(x)\)가 저장된다.

\section{연산자}

\begin{center}
\begin{tabular}{ll}
\toprule
연산자 & 의미 \\
\midrule
\texttt{+} & 덧셈 \\
\texttt{-} & 뺄셈 \\
\texttt{*} & 곱셈 \\
\texttt{/} & 나눗셈 \\
\texttt{//} & 바닥 나눗셈 \\
\texttt{\%} & 나머지 \\
\texttt{**} & 거듭제곱 \\
\bottomrule
\end{tabular}
\end{center}

\begin{lstlisting}
print(10 / 3)
print(10 // 3)
print(10 % 3)
print(10 ** 2)
\end{lstlisting}

Python의 \texttt{//}와 \texttt{\%}는 다음 관계를 만족한다.

\[
a=b(a//b)+(a\%b)
\]

예를 들어

\[
-10=3(-4)+2
\]

이므로

\[
-10//3=-4,\qquad -10\%3=2
\]

이다.

\section{입력과 형 변환}

\begin{lstlisting}
name = input("Name: ")
print(name)
\end{lstlisting}

\texttt{input()}은 항상 문자열을 반환한다. 따라서 숫자로 계산하려면 형 변환이 필요하다.

\begin{lstlisting}
age = int(input("Age: "))
height = float(input("Height: "))
\end{lstlisting}

수학적으로는 다음과 같은 함수 합성으로 생각할 수 있다.

\[
K
\xrightarrow{\ \texttt{input}\ }
\mathrm{String}
\xrightarrow{\ \texttt{int}\ }
\mathrm{Integer}
\]

\section{미니 프로젝트: BMI 계산기}

\[
\mathrm{BMI}
=
\frac{\mathrm{weight}}{\mathrm{height}^{2}}
\]

\begin{lstlisting}
height = float(input("Enter your height (m): "))
weight = float(input("Enter your weight (kg): "))

bmi = weight / (height ** 2)

print(bmi)
\end{lstlisting}

이 프로그램은 변수, 입력, 형 변환, 연산, 출력을 모두 사용하는 첫 번째 완성 프로그램이다.

\section{핵심 정리}

\begin{itemize}
  \item \texttt{print()}는 결과를 화면에 출력한다.
  \item 변수는 객체를 가리키는 이름이다.
  \item 대입은 이름이 가리키는 대상을 변경한다.
  \item Mutable 객체는 내부 상태를 수정할 수 있다.
  \item Immutable 객체는 내부 상태를 수정할 수 없다.
  \item 부동소수점 수는 대부분 정확히 저장되지 않는다.
  \item \texttt{input()}은 항상 문자열을 반환한다.
  \item 숫자 계산에는 \texttt{int()} 또는 \texttt{float()}를 사용한 형 변환이 필요하다.
\end{itemize}

\section{연습 문제}

\begin{enumerate}
  \item \texttt{print()}와 계산의 차이를 설명하시오.
  \item 변수를 객체를 가리키는 이름으로 이해하는 것이 더 정확한 이유를 설명하시오.
  \item Mutable과 Immutable의 차이를 설명하시오.
  \item \texttt{0.1 + 0.2}가 정확히 \texttt{0.3}이 되지 않는 이유를 설명하시오.
  \item \texttt{input()}이 항상 문자열을 반환하는 이유를 설명하시오.
  \item BMI 계산 결과를 소수 둘째 자리까지 출력하도록 코드를 수정하시오.
\end{enumerate}

\end{document}
